% ============================================================================
%  demo.tex — a LaTeX mechanics demonstration and compile smoke test.
%  Neutral content; it exercises the preamble's features (math, units,
%  figures, tables, citations, cross-references) so you can confirm the
%  toolchain works and see how the pieces wire together. NOT a document
%  structure to copy — the sections here are named after LaTeX features.
%  Ships with a local references.bib so it compiles on its own.
%  Compile:  bash scripts/build.sh examples/demo.tex
% ============================================================================
\documentclass[11pt,a4paper]{article}


\input{../assets/preamble}

\title{LaTeX Mechanics Demonstration}
\author{}
\date{}

\begin{document}
\maketitle

\section{Mathematics}
Inline math such as \(a^2 + b^2 = c^2\), and a labelled display equation,
%
\begin{equation}
  \int_{0}^{1} x^2 \, \mathrm{d}x = \frac{1}{3}, \label{eq:demo}
\end{equation}
%
referenced with \cref{eq:demo}. Bold vectors use \(\vec{u}\); paired delimiters
give \(\abs{x}\) and \(\norm{\vec{v}}\).

\section{Units and numbers}
Use \texttt{siunitx} rather than typing units: \SI{22}{\meter\per\second},
\num{1.5e6}, \ang{8}, and \SI{0.5}{\percent}. Inequalities in text need math or
text commands, e.g.\ \(x < 5\) or \textgreater\ 100 (a bare \verb|<| prints
garbage under T1).

\section{Figures and tables}
\Cref{fig:demo} is a vector PDF included at \verb|0.7\textwidth|; \cref{tab:demo}
uses \texttt{booktabs}.

\begin{figure}[H]
  \centering
  \includegraphics[width=0.7\textwidth]{figures/figure.pdf}
  \caption{An example figure.}\label{fig:demo}
\end{figure}

\begin{table}[H]
  \centering
  \caption{An example table.}\label{tab:demo}
  \begin{tabular}{lcc}
    \toprule
    Quantity & Symbol & Value \\
    \midrule
    Alpha & \(\alpha\) & \num{1.0} \\
    Beta  & \(\beta\)  & \num{2.5} \\
    \bottomrule
  \end{tabular}
\end{table}

\section{Citations}
A citation renders as \citep{example-article}; see also \citet{example-book}.

\bibliographystyle{plainnat}
\bibliography{references}

\end{document}
